\chapter{总结与未来展望}
\section{研究总结}
本文以PW对连接成对搬运与腔室周期。两阶段MILP先最小化产品完工时间，再固定完工上界和主要结构改进时序。配置E的30组对照均保持工期并降低等待。配置D的单产品链式工况将总等待由642降至54~s。权重变化使16片实例以16~s额外等待换取8~s节拍偏差下降，24片实例保持不变。

四实例168次对照中，两个混合实例在三个训练种子下均从结构化初始化获益。场景S2的平均工期由521降至492~s。六个2片实例共享SPBS起点的414次比较中，完整选割组合相对HEM加SPBS的实例平均配对改善为0.0776\%，两胜四平。根分离下贪心解码优于束搜索，表明收益随解码设置变化。

单产品链式工况的辅助库存从4片增至8片，已证最优工期由309降至303~s，并行配方保持296~s。四片辅助片可通过重复借用完成链式加工与满片清洁。身份、原模型与事件复核验证了2片324~s最优排程和16片920~s可行排程。

\section{交付与路线等待的双目标权衡}
配对等待下界把指标上限转化为结构筛选条件。两片算例的七个认证点含三个非支持点。四片标称实例的配对下界、领域邻域和对内取放改动均提高了三个种子的超体积。八片与十六片由MOEA/D取得最高HV中位数。75片混合批次得到20个经原模型复验的经验非支配点，两个端点的交付时间相差478~s，等待与节拍指标相差817加权秒。

\section{未来展望}
75片经验前沿仍受生产块和离散次序邻域限制。后续可用工期与等待的联合下界逐段排除未知区间，并改变配方比例、库存和清洁阈值，检验大批次的前沿覆盖度。向配置D的100片批次扩展时，还需测量构造、构模、搜索与验解的总耗时。
